\begin{exercises}
  \item 
    对\nearbyexample{ex:SingJordBlock}做检验。
    \begin{answer}
      我们必须检验
      \begin{equation*}
         \begin{mat}[r]
           3  &0  \\
           1  &3
         \end{mat}
         =
         N+3I=PTP^{-1}
         =
         \begin{mat}[r]
           1/2  &1/2  \\
          -1/4 &1/4
         \end{mat}
         \begin{mat}[r]
           2  &-1  \\
           1  &4
         \end{mat}
         \begin{mat}[r]
           1  &-2  \\
           1  &2
         \end{mat}
      \end{equation*}
      这个计算很容易。
    \end{answer}
  \item 
    每个矩阵都是 Jordan 形式。
    写出它的特征多项式与极小多项式。
    \begin{exparts*}
      \partsitem 
        $\begin{mat}[r]
           3  &0  \\
           1  &3        
         \end{mat}$
      \partsitem 
        $\begin{mat}[r]
           -1  &0  \\
            0  &-1 
         \end{mat}$
      \partsitem 
        $\begin{mat}[r]
           2  &0  &0  \\
           1  &2  &0  \\
           0  &0  &-1/2
         \end{mat}$
      \partsitem 
        $\begin{mat}[r]
           3  &0  &0  \\
           1  &3  &0  \\
           0  &1  &3  \\ 
         \end{mat}$
      \partsitem 
        $\begin{mat}[r]
           3  &0  &0  &0  \\
           1  &3  &0  &0  \\
           0  &0  &3  &0  \\
           0  &0  &1  &3 
         \end{mat}$
      \partsitem 
        $\begin{mat}[r]
           4  &0  &0  &0  \\
           1  &4  &0  &0  \\
           0  &0  &-4 &0  \\
           0  &0  &1  &-4
         \end{mat}$
      \partsitem 
        $\begin{mat}[r]
           5  &0  &0  \\
           0  &2  &0  \\
           0  &0  &3  
         \end{mat}$
      \partsitem 
        $\begin{mat}[r]
           5  &0  &0  &0  \\
           0  &2  &0  &0  \\
           0  &0  &2  &0  \\
           0  &0  &0  &3 
         \end{mat}$
      \partsitem 
        $\begin{mat}[r]
           5  &0  &0  &0  \\
           0  &2  &0  &0  \\
           0  &1  &2  &0  \\
           0  &0  &0  &3 
         \end{mat}$
    \end{exparts*}
    \begin{answer}
      \begin{exparts}
        \partsitem 特征多项式为 $c(x)=(x-3)^2$，
          极小多项式与之相同。
        \partsitem 特征多项式为 $c(x)=(x+1)^2$。
          极小多项式为 $m(x)=x+1$。
        \partsitem 特征多项式为 
          $c(x)=(x+(1/2))(x-2)^2$，
          极小多项式与之相同。
        \partsitem 特征多项式为 $c(x)=(x-3)^3$
          极小多项式与之相同。
        \partsitem 特征多项式为 $c(x)=(x-3)^4$。
          极小多项式为 $m(x)=(x-3)^2$。
        \partsitem 特征多项式为 $c(x)=(x+4)^2(x-4)^2$，
          极小多项式与之相同。
        \partsitem 特征多项式为 
          $c(x)=(x-2)^2(x-3)(x-5)$，
          极小多项式为 $m(x)=(x-2)(x-3)(x-5)$。
        \partsitem 特征多项式为 
          $c(x)=(x-2)^2(x-3)(x-5)$，
          极小多项式与之相同。
      \end{exparts}
    \end{answer}
  \recommended \item
    由所给数据求出 Jordan 形式。
    \begin{exparts}
       \partsitem 矩阵 
         \( T \) 是 \( \nbyn{5} \) 的，仅有特征值 $3$。
         各次幂的零度为：
         \( T-3I \) 的零度为二，\( (T-3I)^2 \) 的零度为三，
         \( (T-3I)^3 \) 的零度为四，而 \( (T-3I)^4 \) 的零度
         为五。
       \partsitem 矩阵 \( S \) 是 \( \nbyn{5} \) 的，有两个特征值。
         对特征值 $2$，零度为：
         \( S-2I \) 的零度为二，而 \( (S-2I)^2 \) 的零度为四。
         对特征值 $-1$，零度为：
         \( S+1I \) 的零度为一。
    \end{exparts}
    \begin{answer}
      \begin{exparts}
         \partsitem 变换 $t-3$ 是幂零的
          （即 $\gennullspace{t-3}$ 是整个空间），
          它通过两条链作用在一个链基上：
          $\vec{\beta}_1\mapsto\vec{\beta}_2\mapsto\vec{\beta}_3
            \mapsto\vec{\beta}_4\mapsto\zero$
          以及 $\vec{\beta}_5\mapsto\zero$。
          因此，$t-3$ 可以用这个标准形来表示。
          \begin{equation*}
            N_3=
            \begin{mat}[r]
               0  &0  &0  &0  &0  \\
               1  &0  &0  &0  &0  \\
               0  &1  &0  &0  &0  \\
               0  &0  &1  &0  &0  \\
               0  &0  &0  &0  &0
             \end{mat}
          \end{equation*}
          从而 $T$ 相似于这个标准形矩阵。
          \begin{equation*}
           J_3=N_3+3I=
           \begin{mat}[r]
               3  &0  &0  &0  &0  \\
               1  &3  &0  &0  &0  \\
               0  &1  &3  &0  &0  \\
               0  &0  &1  &3  &0  \\
               0  &0  &0  &0  &3
             \end{mat}
          \end{equation*} 
         \partsitem 变换 $s+1$ 在子空间 $\gennullspace{s+1}$ 上的限制
          是幂零的，它在链基上的作用为 $\vec{\beta}_1\mapsto\zero$。  
          变换 $s-2$ 在子空间 $\gennullspace{s-2}$ 上的限制
          是幂零的，它在链基上的作用为 $\vec{\beta}_2\mapsto\vec{\beta}_3\mapsto\zero$
          与 $\vec{\beta}_4\mapsto\vec{\beta}_5\mapsto\zero$。        
          因此 Jordan 形式如下。
          \begin{equation*}
            \begin{mat}[r]
              -1  &0  &0  &0  &0  \\
               0  &2  &0  &0  &0  \\
               0  &1  &2  &0  &0  \\
               0  &0  &0  &2  &0  \\
               0  &0  &0  &1  &2
             \end{mat} 
          \end{equation*}
      \end{exparts}  
    \end{answer}
  \item 
    为每个例题求出换基矩阵。 
    \begin{exparts*}
      \partsitem \nearbyexample{ex:FirstJordForm}
      \partsitem \nearbyexample{SecJordanForm}
      \partsitem \nearbyexample{ThirdJordanForm}
    \end{exparts*}
    \begin{answer}
      对每一小题，由于基的取法有很多种，所以还可以有许多别的答案。
      当然，检验一个答案是否给出 Jordan 形式的 $PTP^{-1}$
      的计算才是正确与否的裁判。
      \begin{exparts}
        \partsitem 下面是箭头图。
          \begin{equation*}
            \begin{CD}
              \C^3_{\wrt{\stdbasis_3}}      @>t>T>    \C^3_{\wrt{\stdbasis_3}}   \\
                @V\scriptstyle\identity V\scriptstyle PV  
                                    @V\scriptstyle\identity V\scriptstyle PV \\
              \C^3_{\wrt{B}}                 @>t>J>         \C^3_{\wrt{B}}
            \end{CD}
          \end{equation*}
          从左下角移动到左上角的矩阵是这个。 
          \begin{equation*}
            P^{-1}=\bigl(\rep{\identity}{\stdbasis_3,B}\bigr)^{-1}
                  =\rep{\identity}{B,\stdbasis_3}   
                  =\begin{mat}[r]
                     1  &-2   &0 \\
                     1  &0   &1 \\
                    -2  &0   &0
                   \end{mat}
          \end{equation*}
          从右上角移动到右下角的矩阵 $P$ 是 $P^{-1}$ 的逆。
        \partsitem 我们需要这个矩阵及其逆。
          \begin{equation*}
            P^{-1}=
            \begin{mat}[r]
              1  &0  &3  \\
              0  &1  &4  \\
              0  &-2 &0
            \end{mat}
          \end{equation*}
        \partsitem 将这些广义零空间的基拼接起来，
          就可以作为整个空间的基。
          \begin{equation*}
             B_{-1}=\sequence{\colvec[r]{-1\\ 0 \\  0 \\ 1 \\ 0},
                 \colvec[r]{-1\\ 0 \\ -1 \\ 0 \\ 1}}
            \qquad
             B_3=\sequence{\colvec[r]{1 \\ 1 \\ -1 \\ 0 \\ 0},
                 \colvec[r]{0 \\ 0 \\  0 \\-2 \\ 2},
                 \colvec[r]{-1\\-1 \\  1 \\ 2 \\ 0}}
          \end{equation*}
          换基矩阵是这一个矩阵及其逆。
          \begin{equation*}
            P^{-1}=
            \begin{mat}[r]
              -1  &-1  &1  &0  &-1  \\
              0   &0   &1  &0  &-1  \\
              0   &-1  &-1 &0  &1   \\
              1   &0   &0  &-2 &2   \\
              0   &1   &0  &2  &0   \\
            \end{mat}
          \end{equation*}
      \end{exparts}
    \end{answer}
  \recommended \item 
    求每个矩阵的 Jordan 形式与一个 Jordan 基。
    \begin{exparts*}
      \partsitem 
        \(
        \begin{mat}[r]
          -10  &4  \\
          -25  &10
        \end{mat} \)
      \partsitem 
        \(
        \begin{mat}[r]
           5   &-4 \\
           9   &-7
        \end{mat} \)
      \partsitem 
        \(
        \begin{mat}[r]
           4   &0    &0  \\
           2   &1    &3  \\
           5   &0    &4
        \end{mat} \)
      \partsitem 
        \(
        \begin{mat}[r]
           5   &4    &3  \\
          -1   &0    &-3 \\
           1   &-2   &1
        \end{mat} \)
      \partsitem
        \(
        \begin{mat}[r]
           9   &7    &3  \\
          -9   &-7   &-4 \\
           4   &4    &4
        \end{mat} \)
      \partsitem 
        \(
        \begin{mat}[r]
           2   &2    &-1 \\
          -1   &-1   &1  \\
          -1   &-2   &2
        \end{mat} \)
      \partsitem 
        \(
        \begin{mat}[r]
           7   &1    &2   &2 \\
           1   &4    &-1  &-1\\
          -2   &1    &5   &-1\\
           1   &1    &2   &8
        \end{mat} \)
    \end{exparts*}
    \begin{answer}
      一般的做法是将特征多项式 
      $c(x)=(x-\lambda_1)^{p_1}(x-\lambda_2)^{p_2}\cdots $ 
      因式分解，得到特征值 $\lambda_1$、$\lambda_2$ 等。 
      然后，对每个 $\lambda_i$ 我们求一个
      链基，它是变换 $t-\lambda_i$
      限制在 $\gennullspace{t-\lambda_i}$ 上时的作用的链基：
      计算矩阵 $T-\lambda_iI$ 的各次幂并求出
      相应的零空间，直到这些零空间稳定下来
      （不再变化）为止，此时我们就得到了广义零空间。
      这些零空间的维数（零度）告诉我们
      $t-\lambda_i$ 在广义零空间的某个链基上的作用，
      于是我们可以写出次对角线上 $1$ 的分布模式，
      得到 $N_{\lambda_i}$。
      由此矩阵，与 $\lambda_i$ 相关联的 Jordan 块
      $J_{\lambda_i}$ 立即可得：$J_{\lambda_i}=N_{\lambda_i}+\lambda_iI$。
      最后，在对每个特征值都这样做之后，
      我们把它们拼成标准形。
      \begin{exparts}
        \partsitem 这个矩阵的特征多项式
           是 $c(x)=(-10-x)(10-x)+100=x^2$， 
           所以它只有唯一的特征值 $\lambda=0$。
           \begin{center}
             \renewcommand{\arraystretch}{1.25}
             \begin{tabular}{r|ccc}
                \multicolumn{1}{c}{\textit{幂}~$p$} 
                    &$(T+0\cdot I)^p$ &$\nullspace{(t-0)^p}$
                    &\textit{零度}                   \\ 
                \hline
                $1$  
                &\matrixvenlarge{\begin{mat}[r]
                  -10  &4  \\ 
                  -25  &10
                \end{mat}}
                &$\set{\matrixvenlarge{\colvec{2y/5 \\ y}}
                                     \suchthat
                                     y\in\C}$
                &$1$                                       \\
                $2$  
                &\matrixvenlarge{\begin{mat}[r]
                    0  &0  \\ 
                    0  &0
                \end{mat}}
                &$\C^2$
                &$2$
             \end{tabular}
           \end{center}
           （因此，这个变换是幂零的：
           $\gennullspace{t-0}$ 是整个空间）。
           由这些零度我们知道 $t$ 在某个链基上的作用是
           $\vec{\beta}_1\mapsto\vec{\beta}_2\mapsto\zero$。
           这是 $t-0$ 在 $\gennullspace{t-0}=\C^2$ 上的作用的
           标准形矩阵
           \begin{equation*}
             N_0=
             \begin{mat}[r]
               0  &0  \\
               1  &0
            \end{mat}
           \end{equation*}
           而这就是该矩阵的 Jordan 形式。
           \begin{equation*}
             J_0=N_0+0\cdot I=
             \begin{mat}[r]
               0  &0  \\
               1  &0
            \end{mat}
           \end{equation*}
           注意，若一个矩阵是幂零的，则它的标准形
           等于它的 Jordan 形式。

           我们可以用前一小节的方法找到这样一个链基。
           \begin{equation*}
                 B=\sequence{\colvec[r]{1 \\ 0},
                             \colvec[r]{-10 \\ -25}}
           \end{equation*}
           我们取第一个基向量使得它在 $t^2$ 的零空间中，
           但不在 $t$ 的零空间中。
           第二个基向量是第一个基向量在 $t$ 下的像。
        \partsitem 这个矩阵的特征多项式
           是 \( c(x)=(x+1)^2 \)，所以它是单特征值矩阵。
           （即 $t+1$ 的广义零空间是整个空间。） 
           我们有
           \begin{equation*}
             \nullspace{t+1}=\set{\colvec{2y/3 \\ y}\suchthat
                                     y\in\C} 
             \qquad
             \nullspace{(t+1)^2}=\C^2 
           \end{equation*}
           于是 $t+1$ 在某个相应的链基上的作用是
           $\vec{\beta}_1\mapsto\vec{\beta}_2\mapsto\zero$。
           因此， 
           \begin{equation*}
             N_{-1}
             =
             \begin{mat}[r]
                0  &0  \\
                1  &0 
             \end{mat}
           \end{equation*}
           $T$ 的 Jordan 形式是
           \begin{equation*}
             J_{-1}=N_{-1}+-1\cdot I
             =
             \begin{mat}[r]
               -1  &0  \\
                1  &-1
             \end{mat}
           \end{equation*}
           而从上面的零空间中取向量就得到
           这个链基（也可以有别的取法）。
           \begin{equation*}
             B=\sequence{\colvec[r]{1 \\ 0},
                         \colvec[r]{6 \\ 9}}
           \end{equation*}
        \partsitem 特征多项式
            \( c(x)=(1-x)(4-x)^2=-1\cdot (x-1)(x-4)^2 \) 有两个根，
            它们就是特征值 $\lambda_1=1$ 与 $\lambda_2=4$。

            我们分别处理这两个特征值。
            对 $\lambda_1$，计算 $T-1I$ 的各次幂
            得到
            \begin{equation*}
              \nullspace{t-1}=\set{\colvec{0 \\ y \\ 0}
                                      \suchthat y\in\C}
            \end{equation*}
            而 $(t-1)^2$ 的零空间与之相同。
            于是这个集合就是广义零空间 
            $\gennullspace{t-1}$。
            这些零度表明 $t-1$ 限制在广义零空间上后
            在某个链基上的作用
            是 $\vec{\beta}_1\mapsto\zero$。

            对 $\lambda_2=4$ 作类似的计算，得到这些零空间。
            \begin{equation*}
              \nullspace{t-4}=\set{\colvec{0 \\ z \\ z}
                                      \suchthat z\in\C}
              \qquad
              \nullspace{(t-4)^2}=\set{\colvec{z-y \\ y \\ z}
                                          \suchthat y,z\in\C}
            \end{equation*}
            （$(t-4)^3$ 的零空间与之相同，这也必然如此，因为
            特征多项式中与 $\lambda_2=4$ 相关联的项的
            幂是二，所以
            $t-2$ 限制在广义零空间 $\gennullspace{t-2}$ 上
            是幂零的且指数至多为二\Dash 零空间至多经过
            两次 $t-2$ 的作用就稳定下来。）
            零度上升的模式告诉我们
             $t-4$ 在 $\gennullspace{t-4}$ 的某个相应链基上的作用
            是
            $\vec{\beta}_2\mapsto\vec{\beta}_3\mapsto\zero$。

            把这两个特征值的信息
            合在一起，就得到变换 $t$ 的 Jordan 形式。
            \begin{equation*}
              \begin{mat}[r]
                1  &0  &0  \\
                0  &4  &0  \\
                0  &1  &4
              \end{mat}
            \end{equation*}
            我们可以取零空间中的元素来得到一个合适的基。
            \begin{equation*}
              B=\cat{B_{1}}{B_4}=
               \sequence{\colvec[r]{0 \\ 1 \\ 0},
                          \colvec[r]{1 \\ 0 \\ 1},
                          \colvec[r]{0 \\ 5 \\ 5}}
            \end{equation*}
        \partsitem 特征多项式是
            \( c(x)=(-2-x)(4-x)^2=-1\cdot (x+2)(x-4)^2 \)。

            对特征值 $\lambda_{-2}$
            （这里的下标是 \( -2\)），计算
            $T+2I$ 的各次幂得到这个。
            \begin{equation*}
              \nullspace{t+2}=\set{\colvec{z \\ z \\ z}
                                      \suchthat z\in\C}
            \end{equation*}
            $(t+2)^2$ 的零空间与之相同，所以 
            这就是广义零空间 $\gennullspace{t+2}$。
            于是 $t+2$ 限制在 
            $\gennullspace{t+2}$ 上后，在某个相应的
            链基上的作用是 $\vec{\beta}_1\mapsto\zero$。

            对 $\lambda_2=4$，
            计算 $T-4I$ 的各次幂得到
            \begin{equation*}
              \nullspace{t-4}=\set{\colvec{z \\ -z \\ z}
                                      \suchthat z\in\C} 
              \qquad
              \nullspace{(t-4)^2}=\set{\colvec{x \\ -z \\ z}
                                           \suchthat x,z\in\C}
            \end{equation*}
            于是 $t-4$ 在 $\gennullspace{t-4}$ 的某个
            链基上的作用是
            $\vec{\beta}_2\mapsto\vec{\beta}_3\mapsto\zero$。

